\clearpage\thispagestyle{empty}
\pdfbookmark[0]{Quellen zum Ceres-Bericht vom November 1802}{ceres180211-sources}
\begin{center}\textbf{Quellen zum Ceres-Bericht vom November 1802}\end{center}
\noindent Diese KI-Übertragung in heutiges Deutsch folgt dem Auszug aus Franz Xaver von Zachs Ceres-Bericht vom November 1802 in \href{https://gdz.sub.uni-goettingen.de/id/PPN236007467}{Gauß’ \textit{Werken}, Band~VI} (Göttingen, 1874), S.~227--230. Von Zach berichtet über Gauß’ Rechnungen und zitiert seine abschließende Bemerkung. Die Herausgeber der \textit{Werke} fügen einen Auszug aus einem früheren Bericht über Barnaba Orianis Rechnungen in eckigen Klammern ein.
\par\smallskip\noindent Der \href{https://api.digitale-sammlungen.de/iiif/image/v2/bsb10538598_00506/full/1450,/0/default.jpg}{Erstdruck vom November} steht als Beitrag LI in der \textit{Monatlichen Correspondenz}, Band~VI, S.~492--498. Die Fassung der \textit{Werke} beruht auf \href{https://api.digitale-sammlungen.de/iiif/image/v2/bsb10538598_00509/full/1450,/0/default.jpg}{S.~495}--\href{https://api.digitale-sammlungen.de/iiif/image/v2/bsb10538598_00512/full/1450,/0/default.jpg}{498} des Novemberbeitrags und auf \href{https://api.digitale-sammlungen.de/iiif/image/v2/bsb10538598_00073/full/1450,/0/default.jpg}{S.~67}, \href{https://api.digitale-sammlungen.de/iiif/image/v2/bsb10538598_00074/full/1450,/0/default.jpg}{68} und \href{https://api.digitale-sammlungen.de/iiif/image/v2/bsb10538598_00076/full/1450,/0/default.jpg}{70} des Julibeitrags. Die Gleichungen für Radiusvektor und Breite stehen im Novemberheft auf \href{https://api.digitale-sammlungen.de/iiif/image/v2/bsb10538598_00510/full/1450,/0/default.jpg}{S.~496}.
\par\smallskip\noindent Eine Fußnote auf \href{https://api.digitale-sammlungen.de/iiif/image/v2/bsb10538598_00509/full/1450,/0/default.jpg}{S.~495} berichtigt einen Schreibfehler im Oktoberheft (S.~387). Gemeint ist eine Abnahme der Exzentrizität. Das Novemberheft druckt \(-230''00\) als ersten Koeffizienten der Länge und \(0''38\) ohne Vorzeichen für die jährliche Abnahme der Neigung. In den \textit{Werken} stehen \(-231''00\) und \(+0''38\). Die Übertragung folgt der Lesart der \textit{Werke}.
\par\smallskip\noindent Übertragung und Satz wurden mit KI erstellt. Die Digitalisate der Zeitschrift stammen von der Bayerischen Staatsbibliothek.
\clearpage
